\section{Algebraic hulls and coherent formal fields}
\label{hull:section}

Starting from the contracting state towers, we construct affine families
of conormal covectors with coherent full ambient Taylor coefficients
along both actual foliations. These fields will define the formal
symmetry equations in Section~\ref{sym:section}.
All assertions concerning typical points are with respect to Liouville
volume. We refine the regular set countably many times, including by the
conditional full-measure properties on strong leaves. A statement about a
formal field always concerns its full ambient Taylor series, not just its
restriction to the tangent space of a plaque.

A formal rank-$d$ subbundle is represented by $d$ formal columns
having a nonzero $d$-minor at the center, modulo invertible formal
changes of frame. An affine subbundle is represented by a linear
subbundle after adding a scalar slot and intersecting that slot with
$1$. To specify holonomicity precisely, use fixed ambient coordinates
and a Grassmann coordinate chart for the subbundle. Let $B_\beta(a)$
be its derivative coefficient of multi-order $\beta$ at a point $a$
on an axis. Along a smooth axis path $a(t)$, coherence means
\[
 \frac{d}{dt}B_\beta(a(t))
 =\sum_i \dot a_i(t) B_{\beta+e_i}(a(t)).
\]
Here the coefficients are derivatives rather than factorial-divided
Taylor coefficients. This condition is independent of the ambient
and Grassmann charts by the chain rule. Differentiation on the left is
actual differentiation along the leaf. The higher coefficients on the
right retain all their ambient indices; their existence does not
assert transverse differentiability of an ambient field.

\begin{theorem}[The two formal affine fields]\label{hull:main}
At every point of a conull invariant set there are formal affine fields
$\mathcal A^u,\mathcal A^s$ of horizontal $m$-covectors with the following
properties.
\begin{enumerate}[label=\textup{(\roman*)}]
\item Their values contain $u$ and $w$, respectively. At the center,
the exterior product of any element of $\mathcal A^u$ with any element of
$\mathcal A^s$ equals the fixed nonzero exterior product $u\wedge w$.
\item The fields are equivariant under $f$, changes of section, and flip,
with the same constant and height multipliers as $u,w$. In a flow box
their height dependences are $e^{-h\tau}$ and $e^{h\tau}$.
\item On each entire typical strong stable or strong unstable leaf,
every ambient Taylor coefficient has a smooth representative agreeing
almost everywhere with the regular construction. These representatives
are holonomic: differentiating along the leaf gives the corresponding
higher ambient coefficients. The analogous assertion holds on typical
weak leaves, using flow transport.
\item At every finite order the local representatives, their bounds,
and positive radii can be chosen measurably. On compact plaque subboxes
one may restrict to positive-measure families with uniform finite-order
bounds. No common convergence radius for all orders is asserted.
\end{enumerate}
\end{theorem}

We prove the theorem first for the unstable construction. The stable
construction is obtained by reversing time, or equivalently by flip.
The raw states, normal coordinates, and positive weights are those of
Section~\ref{nf:section}. Write $\mathcal E_j(v)$ for the raw state space
through order $j$, $\Psi_{v,j}$ for its exact normal-coordinate map,
and $N_{v,j}$ for the normal dynamics. All projections commute with these
maps. The order-zero projection is the base label $a$ on the stable axis.
The actual states form a continuous graph over $a$.

\subsection{Polynomial changes of center}

\begin{lemma}\label{hull:recenter}
Fix a finite state level. At sufficiently nearby regular centers on the
same strong stable leaf, the changes between normal state coordinates
are restrictions of polynomial automorphisms. They respect all lower
state projections and carry corresponding actual states into one another.
The same conclusion holds for regular centers related by a fixed flow
translation, using the induced change of section.
\end{lemma}

\begin{proof}
Let $v'$ lie on the stable axis of $v$. In raw coordinates the change
of center is a smooth graph-and-density transformation $C$; both charts
use the same stable plaque. Its conjugate by $\Psi_{v,j},\Psi_{v',j}$ is
a smooth local map $Q$ near the state representing the plaque through
$v'$. Take $v'$ sufficiently close that this state is in a
forward-admissible normal domain at $v$. It is the zero state in the
chart at $v'$. Shrinking a neighborhood of it therefore places the
neighborhood in forward-admissible domains for both descriptions.

Let $Q_n$ be the corresponding change at $f^n v,f^n v'$. Exact
functoriality of graph transformation gives, on the chosen neighborhood,
\begin{equation}\label{hull:transition-identity}
 Q=(N_{v',j}^{(n)})^{-1}\circ Q_n\circ N_{v,j}^{(n)}.
\end{equation}
At the future actual state, the raw chart changes have bounded smooth
norms of each prescribed finite order. The normal changes and their
inverses have tempered norms and radii along each of the two regular
orbits. The arguments of $Q_n$ in
\eqref{hull:transition-identity}, minus its future actual center, decrease
exponentially on a fixed smaller initial neighborhood. Consequently
they eventually fit every prescribed finite-order tempered domain.

For a fixed integer $b$, replace $Q_n$ by its Taylor polynomial of degree
$b$ at the future actual center. Its error after substitution is
$O(e^{-\kappa(b+1)n+\epsilon n})$, for a fixed $\kappa>0$ and arbitrarily
small $\epsilon>0$. On the bounded range in question, postcomposition by
$(N_{v',j}^{(n)})^{-1}$ costs at most $e^{Cn+\epsilon n}$, where $C$
depends on the state level but not on $b$. Choose
$\kappa(b+1)>C+1$. The resulting polynomial approximations converge
uniformly to $Q$.

Their degrees are bounded independently of $n$: at a fixed state level
both iterated normal maps and their inverses have bounded weighted, hence
ordinary, polynomial degree, and $b$ has been fixed. The vector space of
polynomials of this bounded degree is closed under locally uniform
convergence on an open set. Thus $Q$ is polynomial. Applying the same
argument to the inverse change gives a polynomial inverse. The two
polynomial compositions equal the identity on an open set and therefore
identically. Projection compatibility follows first for raw changes and
then for the normal changes.

For a fixed flow shift, use its induced raw graph transformation. It
commutes with $f$, fixes the corresponding central states, and has fixed
finite-order smooth bounds. The same argument applies; closeness of the
two phase-space centers under iteration is unnecessary in this case.
\end{proof}

\subsection{The algebraic tower}

Let $V_j(v)$ be the real Zariski closure of the germ at zero of the actual
states in normal coordinates. Equivalently one may take the complex
Zariski closure of the real evaluations and subsequently use its real
structure. Noetherianity implies that a sufficiently small actual-state
neighborhood already has this closure. Normal dynamics identifies
$V_j(v)$ with $V_j(fv)$.

For a finite cutoff $D$, let $\mathbf e_{j,D}(z)$ be the vector of all
monomials of weight at most $D$, including the constant monomial. Their
evaluation span on the actual germ will be denoted $S_{j,D}(v)$.
Weighted substitution by $N_{v,j}$ is an invertible linear cocycle on
this array. On each weight-$d$ block its sole Lyapunov exponent is $-d$.
The invariant-subspace splitting of Lemma~\ref{nf:blocks} therefore gives
\begin{equation}\label{hull:span-splitting}
 S_{j,D}(v)=\bigoplus_{d\le D} S_{j,D}(v)_d
\end{equation}
at typical centers. The annihilator of this span is the space of
polynomial relations in the array; hence the ideal of $V_j(v)$ is
weighted homogeneous.

\begin{lemma}\label{hull:regular-varieties}
After a conull refinement, all $V_j(v)$ are irreducible and smooth at
zero. Their projections onto lower levels, including the full base,
are submersions at zero. Nearby regular centers on the stable axis
have their variety germs identified by Lemma~\ref{hull:recenter}.
These objects and all finite-dimensional ranks used below are measurable.
\end{lemma}

\begin{proof}
For each finite array, its evaluation span is measurable: evaluate the
continuous actual graph on a countable dense set in successively smaller
base balls and take the stabilized finite-dimensional span. Its
annihilator gives the measurable space $I_D$ of polynomial relations of
ordinary degree at most $D$. In an ambient space $\C^N$, the dimension
$d$ of the closure is the largest size of a coordinate subset $S$ for
which $I_D\cap\C[z_i:i\in S]$ is zero for every $D$. Indeed, a
maximal-dimensional irreducible component has a transcendence-basis
subset among its coordinate generators. Conversely, if every component
has a nonzero relation in the coordinates $S$, the product of those
finitely many relations lies in the full relation ideal. This dimension
test uses only countably many finite-dimensional intersection ranks;
the dimension--transcendence-degree identity is
\citet[Tag~00P0]{StacksProject}.

The spans of the gradients at zero of the spaces $I_D$ stabilize in
$(\C^N)^*$. Choose the first degree attaining their maximal rank and
then a basis by a fixed pivot rule; these are measurable choices.
Every irreducible component contains zero: the connected
group of positive weighted dilations fixes the finite set of components,
and contraction to zero preserves each closed component. Thus the local
dimension at zero is $d$. The tangent-space criterion makes smoothness
at zero equivalent to gradient rank $N-d$; smoothness also forces
irreducibility, since all components meet there
\citep[Tags~00KU, 0B8X, and~00NP]{StacksProject}.
This proves measurability of the central smoothness and irreducibility
condition without selecting components measurably. Projection ranks on
the smooth germs are finite matrix tests. Nonzero minors and sufficiently
small coordinate neighborhoods may then be chosen by countable tests.
The components used below are only existential choices at a fixed
center.

The dimensions of the closures are invariant under $f$ and hence
constant almost everywhere. Work in one such regular stable-axis chart.
Zariski density of the actual graph supplies actual points outside the
singular and intersection loci of the maximal-dimensional components.
It supplies them in arbitrarily small neighborhoods, since the closure
is the closure of a germ. Avoiding these closed algebraic exceptional
loci gives a relatively open subset of the continuous actual graph.
Its regular points have full conditional leaf measure. At any such
point, polynomial re-centering maps its closure into the selected
component. The almost-everywhere equality of dimensions makes this
closure the whole component. Its center is smooth and its closure is
irreducible.

Let $G$ be the measurable invariant set where the central closure is
irreducible and smooth at its center. To pass from the preceding leaf
argument to ambient measure, suppose that $G$ were null. Choose the
regular center above so that its stable plaque meets $G$ in conditional
measure zero, as allowed by the conull convention in the preliminaries.
The relatively open actual-state patch just constructed instead gives
a positive leaf-measure set of regular centers in $G$, a contradiction.
Thus $G$ has positive ambient measure; invariance and ergodicity make
it conull.

For the projection assertion, proceed through the levels.
The projection of the actual germ is the lower actual germ,
so its image is Zariski dense in the lower irreducible variety.
In characteristic zero the dominant projection has full rank on a
nonempty open part of the smooth locus. The same actual-point and
re-centering argument gives a positive leaf-measure patch of centers
with full projection rank. Applying the same conditional-null
contradiction to this invariant rank condition proves the central
submersion assertion almost everywhere.
These are the usual smooth-locus and generic-rank facts for affine
varieties; see \citet{Hartshorne1977} and
\citet[Tags~0B8X, 07ND, and~02K4]{StacksProject}.

The complex formulation introduces no additional component problem.
A selected real actual point outside the other components belongs to
both its component and the complex conjugate of that component. The
two components therefore coincide, so the selected component is defined
over the reals. The selected point is nonsingular; its real locus is
Zariski dense and locally smooth of real dimension equal to the complex
dimension \citep[Theorem~2.2.9(2) and its proof]{Mangolte2019}.
Density of the actual evaluations themselves comes from the closure
argument above. Finally, irreducibility and equal dimension make the
local inclusions supplied by polynomial re-centering equalities of
germs. Intersect the resulting conull sets over the countably many
levels and finite tests.
\end{proof}

We shall always restrict to smooth submersion neighborhoods. In raw
coordinates they have local coordinates $(a,t_j)$, smooth in all
variables and analytic in $t_j$ for fixed $a$, including any fixed
derivatives in $a$. Indeed the inverse normal change has this property
by Lemma~\ref{nf:fiber-analyticity}, and the parameterized implicit function
theorem preserves it.

At a typical center, take the fiber $a=0$ of $V_1(v)$. Its elements give
unstable conormal alternatives at that point, with the prescribed
zero-order density. Higher levels give the same fiber germ at this
level, by submersion. Define stable alternatives in the opposite tower.
These are conormal values supplied by alternative state records. We
will prove their cross-pairing identity first, then construct from their
affine spans full ambient formal fields of constant rank.

\begin{lemma}[Homogeneous parameters]\label{hull:parameters}
The smooth germs $V_j(v)$ have compatible positively weighted
polynomial parametrizations. Independent parameters can be introduced
successively with the state levels. Their new weights tend to infinity
with the level. A finite collection of independent parameters extends
to the full tower by setting all later independent parameters to zero.
If their weights are at most $D_0$, a component of weight $d$ on this
section has ordinary vanishing order at least $d/D_0$.
\end{lemma}

\begin{proof}
The tangent spaces at zero are graded, and their projections are
surjective. Choose graded linear retractions onto them, compatibly with
the projections. The local graph over the tangent space is invariant
under positive weighted dilation. Taylor expansion at zero then makes
each output coordinate a weighted homogeneous polynomial: the Taylor
remainder, pulled toward zero by dilation and divided by the fixed
output weight, tends to zero once the Taylor order is sufficiently high.
Choose successive graded complements for the tangent kernels. Their
weights are among those of the new raw derivative blocks; the minimum
such weight tends to infinity by the raw jet computation. Finally a
monomial of ordinary degree $q$ in parameters of weights at most $D_0$
has weight at most $qD_0$, proving the last assertion.
\end{proof}

\subsection{Extraction from transported actual states}

\begin{lemma}[Pure-weight evaluation]\label{hull:evaluation}
Fix a finite level and monomial array. A vector $b$ in its actual
evaluation span of pure weight $d$ can, along negative returns $-n$,
be written as a linear combination of actual evaluations transported
to the target center. For every $\epsilon>0$ the sum of absolute values
of the coefficients is at most
$C_\epsilon e^{(d+\epsilon)n}$. The transported states tend
exponentially to zero. Arrays can be enlarged before making this choice,
and the larger evaluation spans project surjectively onto the smaller
ones, separately at each weight.
\end{lemma}

\begin{proof}
At each regular source choose an evaluation basis on an arbitrarily
small actual-state neighborhood inside a forward-admissible domain.
Such a basis exists by the definition of the germ span. Countable dense
tests on the actual graph give a measurable choice. Restrict to a
positive-measure set where the domain size, basis norms, and inverse
conditioning are uniformly controlled. Pull $b$ back by the monomial
substitution cocycle. Its norm is at most
$C_\epsilon e^{(d+\epsilon)n}$ by the single-exponent estimate on its
weight block. At negative returns to this set, express it in the chosen
basis and transport the resulting identity forward.

Every basis state is actual, and its transport is an actual state in
the target chart. Forward-admissibility and the positive minimum weight
give exponential convergence to zero. Projection surjectivity follows
because evaluation vectors project to evaluation vectors on the same
actual germ. Equation~\eqref{hull:span-splitting} gives the assertion
weight by weight. Larger cutoffs may alter the conditioning constant
and return set, but not the exponent $d$ of the desired vector.
\end{proof}

Positive-frequency returns will be used throughout. For two towers,
the magnitudes of selected positive and negative return times can be
taken in $[N,2N]$ for all sufficiently large $N$: the ergodic theorem
gives a positive asymptotic count on each such interval. All choices in
Lemma~\ref{hull:evaluation} are made at finite levels, so a countable
regular-set refinement suffices.

\begin{lemma}[Smooth identities and fiber identities]
\label{hull:identities}
If a smooth function vanishes on the local actual states, its Taylor
series at a typical center vanishes on the formal germ of $V_j$.
Moreover, let a function on a local smooth piece of $V_j$ be smooth in
the base and analytic on each fiber. If it vanishes on the nearby
typical actual states, it vanishes on a product neighborhood with
connected fibers. The same assertion applies to each finite coefficient
of a formal calculation having this regularity.
\end{lemma}

\begin{proof}
Extend the function smoothly off $V_j$ if necessary. Fix a weight $d$
in its Taylor series, and take a Taylor polynomial of ordinary degree
$b$ large enough that its remainder on transported states is
$O(e^{-(d+1)n})$. This is possible because all normal variables decrease
at some fixed exponential rate. Include every monomial of this
polynomial in the evaluation array, and apply
Lemma~\ref{hull:evaluation} to a pure weight-$d$ vector. All other
weights are killed exactly. The coefficient cost is
$e^{(d+\epsilon)n}$, whereas the smooth remainder tends to zero faster.
Thus the weight-$d$ Taylor polynomial annihilates the whole evaluation
span and vanishes on $V_j$. Repeat for the finitely many weights needed
at each ordinary order.

For the second assertion first use continuity to include all actual
states on the local graph. Re-center at any nearby typical actual
point using Lemma~\ref{hull:recenter}; its closure is the same variety
germ by Lemma~\ref{hull:regular-varieties}. The first assertion makes
the function flat there in every tangent direction. Its restriction to
the fiber is analytic, hence zero on that connected fiber patch. Typical
base labels are dense and have full conditional measure. Continuity
therefore gives zero throughout a product neighborhood. The argument
can be applied separately to any fixed coefficient.
\end{proof}

\subsection{The pairing of opposite alternatives}

\begin{proposition}\label{hull:pairing}
Every sufficiently small unstable alternative and stable alternative
through the center have the same wedge as the actual pair. This remains
true for their affine spans.
\end{proposition}

\begin{proof}
Use the paired charts and density backgrounds of
Section~\ref{nf:section}. Transport one actual unstable plaque from a
negative return and one actual stable plaque from a positive return.
Write their equations as $x=P(y)$ and $y=P^-(x)$, with log-density
records $r(y),r^-(x)$. At their unique nearby intersection the actual
identity is
\begin{equation}\label{hull:actual-wedge}
 e^{r(y)+r^-(x)}\det\bigl(I-DP(y)DP^-(x)\bigr)=1.
\end{equation}
The two backgrounds cancel at the same point. Base labels and slopes
tend exponentially to zero, so the implicit intersection derivative
is uniformly invertible.

Fix a desired pair of weights $(d,d')$. First truncate spatial graphs
and density records at a large order $R$. Actual plaque derivatives
have uniform bounds of every fixed order. Spatial truncation changes
the intersection and the left side of
\eqref{hull:actual-wedge} by at most
$C_R(|x|+|y|)^{R-1}$; the loss of one derivative accounts for the slopes.
Since the return lengths $n,n'$ are comparable, this is smaller than
$e^{-(d+1)n-(d'+1)n'}$ when $R$ is sufficiently large. The truncated
intersection calculation is a smooth function of finitely many state
coordinates. Taylor-expand it to sufficiently high ordinary order that
the new remainder has the same bound. All low biweight components
stabilize as these orders increase, because the formal intersection
displacements have zero constant term and every state weight is positive.

Take product evaluation arrays containing every term of this finite
polynomial. Pure-weight extraction on the two arrays has coefficient
cost at most $e^{(d+\epsilon)n+(d'+\epsilon)n'}$. The exact identity
on actual plaque pairs and the two remainder bounds force its
$(d,d')$ component to vanish on the product of the two varieties.
This proves the formal identity at every finite order.

Now restrict to the two zero-base fibers. Their intersection equations
have the solution $x=y=0$ identically, so only the two slopes and
prescribed zero-order densities remain. Arbitrary low-level fiber germs
lift to every higher finite level used in the test. Fiber analyticity
upgrades the formal identity to an identity near the actual pair.
Finally bilinearity gives the same constant on affine combinations.
\end{proof}

\subsection{Transport in the opposite direction}

The centered alternatives now satisfy the wedge identity. To turn their
affine spans into formal fields, we must also move the construction along
the central unstable plaque. Polynomial re-centering on the stable axis
does not provide this step: the normal coordinates are only measurable
in their centers.

We will construct a graph-and-density family
$\Phi_v(z,y)=(P_z(y),r_z(y))$. Here $y$ is an actual coordinate on a
small central unstable patch, while $z$ denotes formal normal-state
parameters. The coefficient at parameter zero is the actual central
plaque record. First we construct smooth spatial coefficients with the
prescribed state jets. We then compare their finite weight truncations
with transported actual plaques, to arbitrary prescribed exponential
accuracy. That comparison will let us test the polynomial relations at
a displaced center using actual evaluations.

Denote the positive minimum state weight by $w_0$. A weight cutoff $W$
leaves a finite-dimensional algebra of parameter polynomials, with
products of weight greater than $W$ set to zero. Its positive ideal is
nilpotent of index at most $1+\lfloor W/w_0\rfloor$. A ``parameter
series with smooth spatial coefficients'' means a compatible family of
these finite weight truncations, without a convergence assertion in the
parameters. On a homogeneous tower section from
Lemma~\ref{hull:parameters}, each finite ordinary parameter jet involves
only finitely many weights and hence finitely many levels.

\begin{lemma}[Finite coefficient graph calculus]\label{hull:calculus}
Fix a weight cutoff $W$ and a target spatial derivative order $k$.
There exist integers $r=r(W,k)$ and a constant $C=C(W,k)$ with the
following property. Pushing a weight-truncated graph and log-density
variation from time $-l$ to the target along the central unstable graph
requires only source spatial derivatives through order $r$ at the
contracted central preimage. Comparing two such inputs of tempered
finite coefficient norm costs at most $e^{Cl+\epsilon l}$ in target
$C^k$ norm. The constants $r,C$ are independent of the spatial Taylor
degree used to specify the inputs.
\end{lemma}

\begin{proof}
Put $q=\lfloor W/w_0\rfloor$. The parameter-zero term is the actual
central graph, kept exactly. In the graph-parameter equation solve the
constant term first. Each successive parameter degree is found by
inverting the same central graph Jacobian and substituting already
known lower degrees. At degree at most $q$, Taylor substitution uses at
most $q$ derivatives in the source spatial variable. Taking $k$ target
derivatives adds finitely many derivatives; the density Jacobian adds
one. For example, any $r\ge k+2q+4$ suffices if all graph, inverse,
and density operations are differentiated together.

The composed section map, its height function, and the inverse central
graph parametrization have derivatives through every fixed order bounded
by $e^{A_r l}$ for some $A_r$. This follows by induction from bounded
one-step derivatives and the chain rule; the number of derivatives is
fixed before $l$ varies. Substitution of $N^{-l}z$ on the weight-$W$
array has norm at most $e^{(W+\epsilon)l}$, by its triangular
single-weight estimates. Every coefficient produced by the preceding
nilpotent calculation is a fixed finite sum of products of these
quantities, input derivatives, and central inverse Jacobians.
The number of factors is bounded in terms of $W,k$. Differences are
estimated by replacing one factor at a time. Input bounds are tempered,
so, after choosing smaller exponential slack for their finitely many
factors, all these costs are bounded by $e^{Cl+\epsilon l}$.
The spatial Taylor degree does not enter the number of parameter
operations or the required derivative order $r$.
\end{proof}

\begin{lemma}[Smooth spatial profiles]\label{hull:profiles}
There is a parameter-formal family $\Phi_v(z,y)$ of graph and
log-density records, with smooth coefficients for $y$ in a fixed small
patch of the central unstable axis. Its Taylor series at $y=0$ is the
prescribed tower of raw states. It is exactly equivariant, coefficient
by coefficient, for the section dynamics and $z\mapsto N_vz$.
For every fixed weight cutoff and spatial order its coefficient norms
are tempered along the central orbit.
\end{lemma}

\begin{proof}
Fix $W$. Work first in the ambient normal-state tower and restrict to
the varieties afterward. By Corollary~\ref{nf:finite-weight}, choose
a level $j_0$ after which every new kernel weight exceeds $W$.
For a fixed $j\ge j_0$, let $\pi:\mathcal E_j\to\mathcal E_{j_0}$
be the projection and let $\Pi_j,\Pi_{j_0}$ be the spectral projections
onto weights at most $W$. If $\mathcal R_j$ is a bounded raw right inverse of
$\pi$, then
\[
 \pi\Pi_j\mathcal R_j=\Pi_{j_0}\pi\mathcal R_j=\Pi_{j_0}.
\]
Since $\ker\pi$ has no weight at most $W$, the restriction of
$\Pi_j\mathcal R_j$ to $\operatorname{im}\Pi_{j_0}$ is the inverse of the
low-weight projection. Its norm is tempered. Thus the stabilized
low-weight inputs supply arbitrarily high spatial jets with tempered
bounds at each fixed higher level, without a bound on a measurable
choice of parameters for $V_j$.

At time $-l$, take the spatial Taylor polynomial of degree $R$ of the
prescribed variation, modulo weights greater than $W$, and push it to
the target, using $N^{-l}z$ as source parameter. Keep the parameter-zero
central graph and its density exactly, rather than Taylor-truncating
them. Denote the resulting approximation by $\Phi^{R}_{v,l}$.
The preimage of a fixed small target central patch has spatial size
at most $C e^{-\gamma l}$, for a fixed $\gamma>0$.

Compare consecutive approximations at source time $-l$, first applying
one step to the earlier source polynomial. Exact raw state dynamics
makes their spatial Taylor jets agree through degree $R$. In particular,
the crossing density transformation does not require the next shape
derivative. All higher derivatives needed for this comparison are
tempered for fixed $R,W$. Taylor's formula at the contracted source
argument and Lemma~\ref{hull:calculus} give
\begin{equation}\label{hull:profile-difference}
 \|\Phi^R_{v,l+1}-\Phi^R_{v,l}\|_{C^k}
 \le C_{R,W,k,\epsilon}(v)
 e^{-\{\gamma(R+1-r(W,k))-C(W,k)-\epsilon\}l}.
\end{equation}
Choose $R$ so that the exponent in braces is positive. The differences
are summable in $C^k$. Comparing two spatial cutoffs whose jets agree
through the smaller cutoff gives the same estimate. Thus increasing
$R$ after obtaining $C^0$ convergence does not change the limit.
One may then increase $R$ again to obtain convergence in each $C^k$.
The limits for different $W$ are compatible by the same uniqueness
estimate.

Every prescribed finite spatial jet is retained by taking $R$
sufficiently high; convergence in that derivative order proves Taylor
matching. Shifting the source sequence by one step proves exact formal
equivariance. Finally, all estimates in
\eqref{hull:profile-difference} hold at shifted regular centers with
tempered prefactors. Summing the geometric series after reducing the
slack shows that each fixed coefficient norm of the limit is tempered.
Only a fixed weight cutoff and a fixed spatial norm are being bounded
at a time.
\end{proof}

The following quantitative comparison connects this constructed family
to actual plaques. It is the reason formal identities at a displaced
axis point can be tested by actual evaluations.

\begin{lemma}[Approximation of actual transported plaques]
\label{hull:actual-approximation}
Fix a target norm consisting of graph derivatives through order $k+1$
and log-density derivatives through order $k$, and fix $K>0$.
For a sufficiently large weight cutoff $W$ and a sufficiently high
finite state level, the following holds. Start at a negative return
$-n$ with an actual state in a bounded small forward-admissible domain,
and transport its actual plaque to the target. On a fixed smaller
central unstable patch, substituting its target normal coordinates
into $\Phi_v$ truncated through weight $W$ approximates that plaque
in the prescribed norm with error at most $C e^{-Kn}$.
The level, cutoff, and constant may depend on $k,K$.
\end{lemma}

\begin{proof}
Let the source be at time $-n$ and consider only times from
$-\lfloor n/2\rfloor$ to zero. On every fixed finite normal array,
a monomial of weight $d$ in the transported state is bounded by
\begin{equation}\label{hull:half-small}
 C_\epsilon e^{-dn/2+\epsilon n}
\end{equation}
throughout this last half. Here and below the finite arrays are fixed
before taking $n$ large, and slack is divided among their finitely many
coefficients. Slow regularization bounds every fixed tempered
conversion factor uniformly by $C_\epsilon(v)e^{\epsilon n}$ on
this interval.

We first establish a common geometric patch. The actual crossings
remain on the stable axis. At a crossing, the inverse unstable
parameter differential over a block is the transpose of the
contracting stable-axis differential, because $D_xB(x,0)=0$ and
symplecticity gives $D_yB(x,0)=D_xA(x,0)^{-t}$. Choose a fixed block
length making this inverse norm less than $1/4$, after the bounded
chart conversions. Uniform finite derivatives of actual plaques and
of this fixed block map extend the inverse contraction to a fixed
small parameter ball with norm less than $1/2$. During the last half
the crossings tend uniformly to zero, so the same ball is available
for all sufficiently large $n$. Shorter intermediate blocks have
bounded cost. A fixed small final patch therefore has true inverse
arguments of size
\begin{equation}\label{hull:half-arguments}
 |y_{-\lfloor n/2\rfloor}|\le C e^{-\gamma_1 n}
\end{equation}
at halfway, for a fixed $\gamma_1>0$.

Fix the low-order comparison domain before choosing $W$. It consists
of graph jets through $k+2$ and density jets through $k+1$ in a compact
neighborhood of the bounded central jets, with graph Jacobians bounded
away from singularity. The true jets lie there after shrinking the
patch. Let $L_k\ge2$ bound the one-step jet transformation and its
first derivatives on a slightly larger such domain. This number uses
only the fixed low orders and the uniform raw geometric bounds.

For each fixed $W$, every positive-parameter term of the truncated
profile tends uniformly to zero in these low norms during the last
half: its spatial coefficient is tempered, while its state monomial
has positive weight and satisfies \eqref{hull:half-small}. Thus the
substituted profiles lie in the same comparison domain for all
$n\ge n_0(W)$. This is the decisive point: the threshold depends on
$W$, but the amplification constant $L_k$ does not.

Compare the actual and substituted profiles at halfway. Choose a
spatial order $S$ and a finite tower containing their spatial jets
through $S$. Expand the smooth inverse normal change to ordinary
degree $b$, where $(b+1)w_0>2W$. All terms of weight at most $W$
agree by Lemma~\ref{hull:profiles}. Every remaining polynomial term
has weight greater than $W$ and is bounded using
\eqref{hull:half-small}. The ordinary Taylor remainder satisfies
the same bound by the choice of $b$. Consequently the differences
of spatial Taylor coefficients through order $S$ are at most
$C e^{-Wn/4}$, after sufficiently small slack.
The spatial Taylor remainder at the true inverse argument is bounded
by $C e^{-\gamma_1(S-k)n+\epsilon n}$, using
\eqref{hull:half-arguments}, uniform actual plaque derivatives, and
tempered profile coefficients. Choose
\begin{equation}\label{hull:spatial-choice}
 \gamma_1(S-k)>W/2+1.
\end{equation}
The initial error in the graph/density norm at halfway is then at most
$C e^{-Wn/4}$.

At every last-half step, the substituted truncated profiles satisfy
the exact graph/density recurrence modulo weights greater than $W$.
For fixed $W$ and each fixed parameter derivative order, the real
graph transformation of the finite polynomial profile is defined on
a temperedly small parameter ball and has tempered derivatives there.
Indeed its central graph Jacobian is uniformly invertible; the finite
profile coefficient bounds give a tempered radius on which it remains
invertible, and implicit differentiation uses only finitely many
tempered coefficients. Exponentially small transported states
eventually lie in these balls, uniformly over the last-half interval.
Taylor-expand that smooth recurrence to ordinary degree $b$ with
$(b+1)w_0>2W$. Formal equivariance kills its terms through weight
$W$, and \eqref{hull:half-small} bounds all remaining terms and the
smooth remainder by $C e^{-Wn/4}$. The same estimate holds in the
fixed low derivative norms. The constant and the large-$n$ threshold
may depend on $W,S$, which have now been fixed.

These errors can be compared at the true successive arguments.
The approximate transformed argument differs from the true one by
a bounded multiple of the source point error. Moving its jets to the
true argument costs at most another bounded factor, because one
additional spatial derivative lies in the fixed comparison domain.
Increasing $L_k$ once, if necessary, the successive errors $E_i$
satisfy
\[
 E_{i+1}\le L_k E_i+C e^{-Wn/4}.
\]
There are at most $n$ steps. Hence the final error is at most
\begin{equation}\label{hull:actual-error}
 C n L_k^n e^{-Wn/4}
 \le C' e^{-\{W/4-\log L_k-1\}n}.
\end{equation}
Finally choose $W>4(K+\log L_k+2)$, then $S$ by
\eqref{hull:spatial-choice}, then the requisite finite level and
ordinary Taylor degrees. This order of choices proves the result.
All estimates were uniform on the fixed final subpatch.
\end{proof}

\begin{lemma}[Translation of the formal family]\label{hull:transport}
At a nearby typical point of the central unstable plaque, re-center
$\Phi_v(z,y)$ over that point's stable axis. Its graph-and-density
state series belongs formally, at every finite arrival level, to the
arrival variety. The re-centered zero-order log-density has exactly
the prescribed background value. This assertion is understood on
finite homogeneous parameter sections and then in weight completion.
\end{lemma}

\begin{proof}
Fix the arrival chart and a polynomial relation of its finite-level
variety. Re-centering is a smooth graph transformation followed by
an implicit equation for the crossing. At parameter zero it describes
the actual central unstable plaque, and its crossing derivative is
invertible. Substituting $\Phi_v$ therefore gives a well-defined
parameter series with smooth spatial coefficients.

Fix one weight $d$ of this series. Evaluate on actual states transported
from negative returns as in Lemma~\ref{hull:evaluation}. At the arrival
point the transported actual plaques satisfy the chosen polynomial
relation exactly. By Lemma~\ref{hull:actual-approximation}, the
substituted profiles approximate all graph and density jets needed
for the test to accuracy $e^{-Kn}$ for arbitrary prescribed $K$.
The arrival chart is fixed, so re-centering and polynomial evaluation
have fixed smooth comparison constants. Taylor-expand the resulting
composition in positive-weight variables through sufficiently high
ordinary order, then retain weights through a large cutoff. On
transported states its discarded terms again decrease as fast as
required, by the positive minimum weight and
\eqref{hull:half-small} at the final time.

Choose $K>d+1$ and take an evaluation array containing every term of
this finite polynomial. Pure-weight extraction has cost
$e^{(d+\epsilon)n}$, so its weight-$d$ component vanishes on the
departure variety. Increasing the level is harmless by the
surjectivity in Lemma~\ref{hull:evaluation}. The same test, applied to
the re-centered zero-order log-density minus its prescribed background,
proves that equality too; omitting this entry from the state causes no
loss. Repeat at every finite weight and arrival level.

On any fixed homogeneous parameter section, an ordinary coefficient
uses only bounded weights by Lemma~\ref{hull:parameters}. Thus all
compositions in the statement are ordinary formal compositions at
each finite test. No analyticity in the spatial variable, convergence
in parameters, or regular dependence of the arrival normal chart on
the center has been used.
\end{proof}

\subsection{Taking the spans}

The varieties constructed above parameterize alternative plaque jets.
We now use their raw state coefficients to form spatial Taylor series,
and take the affine span of their conormal forms at a fixed formal point.
The remaining task is to prove constant formal rank and independence
from the parameters and higher lifts.

Choose compatible local raw submersion coordinates $(a,t_j)$ on each
finite level. For a fixed Taylor order in $y$, take $j$ large enough to
contain its graph and log-density coefficients. These coefficients are
smooth in the base $a$ and analytic in the hidden
parameters $t_j$, also after any fixed number of base derivatives, by
Lemma~\ref{nf:fiber-analyticity}. Thus the spatial series used here are
formal in $y$, with these actual coefficient functions of $(a,t_j)$.
When we re-center at a typical actual state and Taylor-expand in the
parameters, these are precisely the state jets matched by
Lemma~\ref{hull:profiles}. Its smooth spatial coefficients will let us
transport those parameter Taylor coefficients along the actual unstable
plaque; no profile is being evaluated at arbitrary hidden parameters.

Restore the fixed form units in the graph-conormal formula. To retain
the translation part when passing from affine to linear spans, augment
each alternative with a scalar slot and write
\[
 \mathbf u_{a,t_j}(y)=(u_{a,t_j}(y),1).
\]
At a formal section point $(x,y)$, solve the graph equation for the base
label while leaving the hidden parameters free:
\begin{equation}\label{hull:point-constraint}
 \begin{split}
 P_{b(x,t_j,y),t_j}(y)&=x,\\
 b(x,t_j,0)&=x,\\
 S(x,t_j,y)&=\mathbf u_{b(x,t_j,y),t_j}(y).
 \end{split}
\end{equation}
The submersion to the base gives this formal solution in $y$, with
coefficients smooth in $x$ and analytic in $t_j$. We take the span of
the augmented forms by taking all derivatives of $S$, including order
zero, in the free parameters while keeping $(x,y)$ fixed. At the central
point its dimension is a number $d$, constant almost everywhere by
equivariance and ergodicity. It is already detected at a finite level:
at the center only the slope and prescribed density enter, and the
fiber dependence is analytic.

\begin{lemma}[Constant formal rank]\label{hull:rank}
The constrained augmented span is a free formal subbundle of rank
$d$. Its coefficients can be taken smooth in the stable-axis variable,
with their full coherent ambient jets. The resulting subbundle is
independent of the hidden parametrization and of all higher-level
lifts used in its construction.
\end{lemma}

\begin{proof}
Fix a spatial Taylor order. Each fixed $y$-order in
\eqref{hull:point-constraint} uses only finitely many shape and density
jets, including one additional shape derivative for the conormal slope.
Increasing the number of derivatives in $t_j$ does not increase that
required state level: it differentiates the same functions on that level.

For any $d+1$ multi-indices $\nu_0,\ldots,\nu_d$, including zero,
form the wedge of the columns $\partial_{t_j}^{\nu_i}S$ while keeping
$(x,y)$ fixed. Substitute $x=P_{a,t_j}(y)$ afterward. Every spatial
Taylor coefficient of this wedge is a function on a finite $V_j$,
smooth in $a$ and analytic in $t_j$. We claim that it vanishes.

First center the calculation at a nearby typical actual state.
Polynomial stable-axis re-centering identifies its variety germ
with the present finite-level germ. In the re-centered chart choose
a homogeneous tower section over a level containing all jets needed
for the tested spatial coefficient. This section covers the entire
finite-level parameter germ after the invertible chart comparison.
Higher choices do not affect the tested spatial coefficient at
displacement zero.

Apply Lemma~\ref{hull:profiles} in this re-centered chart. At each
small finite displacement on its actual unstable plaque, the
parameter-zero graph is actual. The derivative of its crossing with
respect to the base parameter is invertible near displacement zero,
so the fixed-point constraint can be solved there as a formal series
in the finite parameters. At typical displaced centers,
Lemma~\ref{hull:transport} places all these constrained variations in
the arrival centered fibers. Their augmented-form derivatives belong
to the centered $d$-dimensional span. Every $(d+1)$-fold wedge
therefore vanishes.

Each such derivative has smooth dependence on spatial displacement.
Typical displaced points are dense on this plaque, so the wedge
vanishes smoothly along a patch. Its tested Taylor coefficient at
zero is precisely the coefficient under consideration at the
chosen actual state: finite-level re-centering preserves the needed
jets, and the homogeneous lift covers their finite parameter germ.
Thus this coefficient vanishes at every nearby requisite typical
actual state. Lemma~\ref{hull:identities} now makes it vanish
throughout a product neighborhood in $V_j$. At a fixed spatial
order all multi-indices can use the same finite level and connected
fiber patch. The conclusion applies to each of their analytic
coefficient functions there.

Choose $d$ independent derivative columns at the central state, and
let $B$ be their formal continuation in the section-point variables.
A corresponding $d$-minor is a unit. Substitute the formal solution
of \eqref{hull:point-constraint} in the wedge identities just proved.
They show that every derivative column lies in the free rank-$d$
module generated by $B$. Differentiating a basis column in a hidden
direction gives another derivative column; hence
\[
 \partial_{t_j}B=B C_j
\]
for a formal coefficient matrix $C_j$. Normalize the selected minor
of $B$ to the identity. Its hidden derivative is then zero, since
the same relation has zero rows on that minor. This proves that the
span is independent of hidden variables.

The chain rule under an invertible hidden reparametrization gives
the same derivative span in both directions. Higher-level parameters
are also hidden parameters, so changing a higher lift cannot alter
the span. For each fixed ambient order choose a sufficiently high
finite level; the unit-minor characterization makes the resulting
jets consistent. Finally, one may retain the base variable $x$ as
an actual smooth variable while expanding only in $y$ in
\eqref{hull:point-constraint}. A smooth reference hidden section
then gives smooth coefficients on the stable axis, whose base
derivatives are exactly their higher ambient coefficients.
\end{proof}

Intersect this augmented linear bundle with scalar slot $1$. It
defines the formal affine field $\mathcal A^u$. The scalar-$1$
intersection is nonempty at the center because it contains the actual
form. The same unit-minor description makes it a formal affine
bundle on the local germ.

\begin{lemma}[Coherence on both axes]\label{hull:coherence}
The full ambient jets of $\mathcal A^u$ have smooth holonomic
representatives along both entire typical strong axes. They agree
with the regular construction almost everywhere and are compatible
with dynamics, changes of section, and regular flow translations.
\end{lemma}

\begin{proof}
Stable-axis agreement at nearby typical centers follows from
Lemma~\ref{hull:recenter} and the intrinsic constrained-span
description. Smooth stable-axis coherence was proved in
Lemma~\ref{hull:rank}.

For the unstable axis take a homogeneous tower section over a level
containing the base variables and the parameters furnishing the
$d$ independent columns. Use $\Phi_v$ to carry the same constrained
basis calculation to a small finite $y$. The columns and every one
of their formal point coefficients depend smoothly on $y$. Their
central minor remains invertible after shrinking the patch.
At a typical arrival center, Lemma~\ref{hull:transport} places their
full graph-state series in every finite arrival variety, including
the prescribed zero-order density. For a prescribed point-jet order,
take an arrival state level containing the required graph and density
derivatives, including the extra graph derivative in the conormal form.
Taylor re-expansion of the smooth spatial profile about this fixed
arrival center is exactly the graph reconstructed from those state
jets. Fix a full formal arrival point
$q=(x',y')$ and a finite jet order. Write $S_{\mathrm{arr}}(q,t')$
for the arrival constrained augmented family from
\eqref{hull:point-constraint}. In an arrival submersion chart, the
state inclusion and the point constraint give, at the chosen finite
order,
\[
 S_{\mathrm{tr}}(q,t)
   =S_{\mathrm{arr}}\bigl(q,\Theta(q,t)\bigr),
\]
where $t$ denotes the departure hidden parameters, $\Theta(q,t)$
the induced arrival hidden parameters, and both augmented families
are expressed in the arrival chart. Only finitely many state levels
enter this identity. Differentiating in $t$ while keeping $q$ fixed
shows by the chain rule that every transported column belongs to
the arrival constrained span. That span also has rank $d$, so the
retained independent columns give the whole bundle. Inclusion and
equality of ranks suffice; no surjectivity between infinite towers
is required. Equality at each finite order proves full ambient
coherence on this unstable-axis patch.

All operations commute with raw graph transformation, normal
comparison, and the appropriate conormal height factors.
The constant height multiplier at a fixed point is independent of
the alternative being spanned. Thus the constructed jets are
equivariant for dynamics and changes of section. Regular centers
related by a fixed flow shift have compatible jets by
Lemma~\ref{hull:recenter}.

To cover an entire typical stable leaf, fix a finite collection
of orders. Local smooth representatives have a positive radius
at almost every center. Choose forward returns on which their
radius is bounded below and the required finite construction
bounds are bounded above. Every compact portion of the central
stable leaf contracts into these patches; pull the representatives
back. Two pulled representatives agree on the dense set of
typical points and hence everywhere. The same argument backward
in time covers the entire unstable leaf. Repeat one finite
collection of orders at a time. Agreement of the orders and their
derivative identities follows on overlaps by density and smoothness.
This does not require a single radius valid at all orders.
Measurability of the radii used here is established below.
\end{proof}

\begin{lemma}[Measurable finite-order bounds]\label{hull:measurability}
The construction and the representatives in
Lemma~\ref{hull:coherence} admit measurable finite-order norm and
radius tests. Smooth flow transport supplies coherent
representatives on typical weak plaques, agreeing leaf-almost
everywhere with the regular data. For almost every individual
strong or weak plaque, every fixed finite collection of coefficient
norms is finite on each compactly contained plaque subpatch.
Positive-measure families of these plaques can therefore be chosen
with a common finite bound on the prescribed subpatches.
\end{lemma}

\begin{proof}
Raw leaf charts and actual finite plaque jets have measurable,
indeed continuous local transverse dependence. The adapted
normal changes have measurable finite jets and domains. All
subsequent finite constructions use polynomial spans, kernels,
nonzero minors, implicit coordinate charts, or limits of measurable
coefficient functions. Countable choices of candidate minors and
coordinate balls give measurable versions. The coefficients of
the resulting formal fields are therefore measurable.

It remains to justify the patch bounds without presupposing a
smooth transverse choice of representative. In a fixed countable
plaque box, the smooth representative on a typical plaque agrees
almost everywhere with these measurable coefficients. Its norm
on a compact subpatch is the essential supremum of the same
coefficient norms there; this is measurable by conditional
integration. Derivative norms are tested in the same way using
the higher coherent coefficients. For a radius bound, work in
a fixed Grassmann chart with a chosen minor bounded away from
zero and impose bounded jets and Taylor difference estimates
through an order strictly higher than the differentiability being
requested, between almost every pair of points on a smaller
plaque ball. These tests are measurable by Fubini and countably
many radii and bounds. Existing smooth coherent representatives
satisfy them locally. Conversely, such bounds extend the tested
jets continuously to the closure of a dense valid set and give
their stipulated finite differentiability and coherence. A dense
countable set of mutually valid points may be selected by
successive Fubini choices. These tests provide the measurable
positive radii required in the globalization proof.

On weak plaques use the whole-strong-leaf representatives and
smooth flow transport. At points for which both constructions
are regular, their jets agree by the flow comparison already
proved. Strong-leaf conditional agreement and flow Fubini then
give agreement almost everywhere on the weak plaque; agreement
at every flow time is unnecessary. Compact subpatches have
finite smooth norms. Conditional essential suprema make these
bounds measurable on the family of weak plaques as well.

An almost-everywhere finite measurable bound is bounded on a
positive-measure subfamily of any positive-measure plaque-label
set. The same applies to the reciprocal of a positive radius
and to any finite list of bounds. Absolute continuity of the
strong and weak foliations and a countable cover allow all
conditional agreement requirements to be imposed simultaneously,
also at integer iterates. These statements concern bounds along
plaques; they do not assert that the formal fields themselves
have tempered norms along every regular orbit.
\end{proof}

\begin{proof}[Proof of Theorem~\ref{hull:main}]
The unstable affine field is supplied by
Lemma~\ref{hull:rank}. Its central value contains the actual
augmented form and hence the actual form. Lemmas~\ref{hull:coherence}
and~\ref{hull:measurability} give the full ambient coherence,
globalization, and measurable finite-order bounds. Restore the
height dependence $e^{-h\tau}$.

Apply the same construction with time reversed to obtain
$\mathcal A^s$, restoring height dependence $e^{h\tau}$. One may
choose charts measurably in flip pairs, with the axes interchanged;
all graph comparisons preserve the construction, so it is compatible
with flip. Proposition~\ref{hull:pairing} gives the constant cross
pairing of the centered affine spans. The transformation laws follow
from the exact raw form transformation and the fixed-point nature of
the span. This proves every assertion of the theorem.
\end{proof}
